\documentclass[aspectratio=169]{beamer}
\usepackage[utf8]{inputenc}
\usepackage[T1]{fontenc}
\usepackage[scaled=0.82]{beramono} % monoespaçada com negrito forte
\usepackage[brazilian]{babel}
\usepackage{graphicx}
\usepackage{array}
\usepackage{tikz}
\usetikzlibrary{positioning,arrows.meta,shapes.geometric,calc,fit}
\usetheme{default}
\usecolortheme{orchid}
\usefonttheme[onlymath]{serif}
\beamertemplatenavigationsymbolsempty
\hypersetup{pdfpagemode=FullScreen}

% ---------------------------------------------------------------------
% Cores (paleta única das figuras)
%   ufpr       contorno externo das figuras, setas, destaques escuros
%   ufprlight  aplicações
%   kfill      camadas de software e hardware (kernels, runtime, hipervisor)
%   kdraw      contorno de todas as caixas internas
%   hostc      processadores e kernel hospedeiro
%   ufpr!20    memória e réplicas de estado
%   ink        texto das figuras
\definecolor{ufpr}{RGB}{52,92,135}
\definecolor{ufprlight}{RGB}{214,226,238}
\definecolor{hostc}{RGB}{200,200,200}
\definecolor{kfill}{HTML}{EFF3F7}\definecolor{kdraw}{HTML}{CBD5E1}
\definecolor{ink}{HTML}{1E293B}
\setbeamercolor{frametitle}{fg=ufpr}
\setbeamercolor{structure}{fg=ufpr}
\setbeamercolor{block title}{bg=ufpr,fg=white}
\setbeamercolor{block body}{bg=ufprlight}
\setbeamercolor{alerted text}{fg=black}
\setbeamertemplate{itemize items}[circle]
% Espaçamento uniforme entre itens em todas as listas
\let\origitemize\itemize
\renewcommand{\itemize}{\origitemize\small\setlength{\itemsep}{3pt}\setlength{\parsep}{0pt}}


% Estilos das figuras
\tikzset{
    card/.style={draw=ufpr, thick, fill=white, rounded corners=3pt},
    chip/.style={font=\ttfamily\tiny\bfseries, text=ink, inner sep=2pt},
    tag/.style={chip, draw=kdraw, fill=kfill, rounded corners=1pt},
    ltitle/.style={font=\scriptsize\bfseries, text=ink},
    head/.style={font=\ttfamily\small\bfseries, text=ink},
    mkcap/.style={font=\scriptsize\itshape, text=gray, align=center},
}
% \layer{xcenter}{ytop}{altura}{preenchimento}{borda}{titulo}
\newcommand{\layer}[6]{%
    \draw[fill=#4, draw=#5, rounded corners=2pt] (#1-19,#2) rectangle (#1+19,#2-#3);
    \node[ltitle] at (#1,#2-3) {#6};}

\newcommand{\refs}[1]{{\scriptsize\color{gray}#1}}

% Numeracao de paginas (centro inferior, entre o logo e o grafismo)
\setbeamertemplate{footline}{%
    \hfill{\usebeamerfont{page number in head/foot}\footnotesize\color{gray}\insertframenumber\,/\,\inserttotalframenumber}\hfill\vspace*{3mm}}


\tikzset{
    dbox/.style={draw=ufpr, fill=ufpr, text=white, rounded corners=2pt, align=center, font=\footnotesize},
    arr/.style={-{Stealth[length=2mm]}, thick, ufpr},
}

% ---------------------------------------------------------------------
% Titulo (template UFPR)
\makeatletter
\newcommand\titlegraphicii[1]{\def\inserttitlegraphicii{#1}}
\titlegraphicii{}
\setbeamertemplate{title page}
{
    \vbox{}
    {\usebeamercolor[fg]{titlegraphic}\inserttitlegraphic\hfill\inserttitlegraphicii\par}
    \begin{centering}
        \begin{beamercolorbox}[sep=2pt,center]{institute}
      \usebeamerfont{institute}\insertinstitute
        \end{beamercolorbox}
        \begin{beamercolorbox}[sep=12pt,center]{title}
      \usebeamerfont{title}\inserttitle\par%
      \ifx\insertsubtitle\@empty%
      \else%
          \vskip0.2em%
              {\usebeamerfont{subtitle}\usebeamercolor[fg]{subtitle}\insertsubtitle\par}%
      \fi%
        \end{beamercolorbox}%
        \vskip0.2em\par
        \begin{beamercolorbox}[sep=2pt,center]{author}
      \usebeamerfont{author}\insertauthor
        \end{beamercolorbox}
        \begin{beamercolorbox}[sep=4pt,center]{date}
      \usebeamerfont{date}\insertdate
        \end{beamercolorbox}
    \end{centering}
}
\makeatother

\title{K-ORCH: orquestrador para instâncias multikernel distribuídas}
\subtitle{Plano de trabalho}
\author{\textbf{Eduardo Kaluf}\\[1pt]{\footnotesize eduardokaluf@ufpr.br}\\[3pt]{\small Orientador: Prof. Carlos A. Maziero}}
\institute{UNIVERSIDADE FEDERAL DO PARANÁ \\
SETOR DE CIÊNCIAS EXATAS \\
DEPARTAMENTO DE INFORMÁTICA}
\date{Novembro de 2026}

\begin{document}

% =====================================================================
    {%
        \usebackgroundtemplate{\centering\includegraphics[width=\paperwidth]{CapaUFPR.png}}
        \setbeamertemplate{footline}{}
        \setbeamertemplate{headline}{\vskip2mm\hfill{\footnotesize\color{gray}\insertframenumber\,/\,\inserttotalframenumber}\hspace*{4mm}}
        \begin{frame}
            \vspace{-12mm}
            \titlepage
        \end{frame}
    }

    {%
        \usebackgroundtemplate{\centering\includegraphics[width=\paperwidth]{FundoUFPR2.png}}

% =====================================================================
        \begin{frame}{Evolução da arquitetura multikernel}
            \centering
            \vspace{4mm}
            \begin{tikzpicture}[x=12.9mm]
    \draw[ufpr, very thick, -{Stealth}] (0,0) -- (10.6,0);
    \foreach \x/\y/\n/\h/\c in {
        0.4/1995/Hive/0.45/ufpr,
        1.5/1997/Disco/1.05/black,
        2.6/2000/{Cellular Disco}/0.45/black,
        3.7/2008/Corey/1.05/black,
        4.8/2009/{Barrelfish}/0.45/ufpr,
        5.9/2012/FusedOS/1.05/black,
        7.0/2014/IHK/0.45/black,
        8.1/2015/Pisces/1.05/black} {
        \fill[ufpr] (\x,0) circle (2.2pt);
        \draw[gray] (\x,0) -- (\x,\h-0.12);
        \node[font=\scriptsize\bfseries, text=\c, above] at (\x,\h-0.15) {\n};
        \node[font=\scriptsize, below] at (\x,-0.12) {\y};
    }
    \fill[ufpr] (9.6,0) circle (3.2pt);
    \draw[gray] (9.6,0) -- (9.6,0.33);
    \node[font=\scriptsize\bfseries, ufpr, above] at (9.6,0.3) {Linux};
    \node[font=\scriptsize, below] at (9.6,-0.12) {2025};
            \end{tikzpicture}

            \vspace{8mm}
            \begin{columns}[T]
                \begin{column}{0.38\textwidth}
                    \textbf{\small Multiprocessadores (1995--2000)}\\[1pt] % exceção: título menor para caber numa linha
                    \setlength{\leftmargini}{1em}\begin{itemize}\item VMs especializadas e contenção de falhas em ``clusters virtuais''\end{itemize}
                \end{column}
                \begin{column}{0.30\textwidth}
                    \textbf{\small Manycore (2008--2009)}\\[1pt] % exceção: título menor para caber numa linha
                    \setlength{\leftmargini}{1em}\begin{itemize}\item Kernel único por núcleo e comunicação explícita por mensagens\end{itemize}
                \end{column}
                \begin{column}{0.32\textwidth}
                    \textbf{\small HPC e Linux (2012--2025)}\\[1pt] % exceção: título menor para caber numa linha
                    \setlength{\leftmargini}{1em}\begin{itemize}\item Kernels leves ao lado do Linux e instâncias Linux particionadas\end{itemize}
                \end{column}
            \end{columns}

        \end{frame}

% =====================================================================
        \begin{frame}{Arquitetura multikernel: Hive}
            \begin{columns}[c]
                \begin{column}{0.50\textwidth}
                    \centering
                    \begin{tikzpicture}[x=1mm,y=1mm]
                        % maquina
                        \draw[ufpr, thick, rounded corners=3pt] (-35,31) rectangle (35,-27);
                        % aplicacoes e imagem unica
                        \draw[fill=ufprlight, draw=kdraw, rounded corners=2pt] (-32,28) rectangle (32,22);
                        \node[font=\scriptsize\bfseries, text=ink] at (0,25) {Aplicações};
                        \draw[fill=kfill, draw=kdraw, rounded corners=2pt] (-32,20) rectangle (32,14.5);
                        \node[font=\scriptsize\bfseries, text=ink] at (0,17.25) {Imagem única do sistema};
                        % celulas
                        \foreach \i/\x in {1/-24,2/-8,3/8,4/24} {
                            \ifnum\i=3
                            \draw[fill=hostc!45, draw=gray, dashed, rounded corners=2pt] (\x-7.3,12) rectangle (\x+7.3,-15.5);
                            \else
                                \draw[fill=white, draw=kdraw, rounded corners=2pt] (\x-7.3,12) rectangle (\x+7.3,-15.5);
                            \fi
                            \node[font=\tiny\bfseries, text=ink] at (\x,9.5) {Célula \i};
                            \draw[fill=kfill, draw=kdraw, rounded corners=1pt] (\x-6,7) rectangle (\x+6,2);
                            \node[font=\tiny\bfseries, align=center] at (\x,4.5) {Kernel Hive};
                            \draw[fill=hostc, draw=kdraw] (\x-6,0) rectangle (\x-0.5,-5);
                            \draw[fill=hostc, draw=kdraw] (\x+0.5,0) rectangle (\x+6,-5);
                            \node[font=\tiny\bfseries] at (\x-3.25,-2.5) {CPU};
                            \node[font=\tiny\bfseries] at (\x+3.25,-2.5) {CPU};
                            \draw[fill=ufpr!20, draw=kdraw] (\x-6,-7) rectangle (\x+6,-12);
                            \node[font=\tiny\bfseries] at (\x,-9.5) {Memória};
                        }
                        % falha
                        \node[font=\tiny\bfseries, text=ink] at (8,-13.8) {falha};
                        % interconexao
                        \draw[fill=kfill, draw=kdraw, rounded corners=1.5pt] (-32,-19) rectangle (32,-24);
                        \node[font=\tiny\bfseries, text=ink] at (0,-21.5) {Interconexão: RPC e memória compartilhada};
                        \foreach \x in {-24,-8,8,24} \draw[ufpr, thick] (\x,-15.5) -- (\x,-19);
                    \end{tikzpicture}
                \end{column}
                \begin{column}{0.50\textwidth}
                    \begin{itemize}
                        \item Divide um multiprocessador em células que executam sua própria instância do kernel Hive
                        \item Falhas de hardware ou software ficam contidas na célula atingida
                        \item Comunicação por memória compartilhada e RPC
                        \item Política global em um processo de usuário
                    \end{itemize}
                \end{column}
            \end{columns}
        \end{frame}

% =====================================================================
        \begin{frame}{Arquitetura multikernel: Barrelfish}
            \begin{columns}[c]
                \begin{column}{0.45\textwidth}
                    \begin{itemize}
                        \item Trata a máquina física como uma rede de núcleos independentes
                        \item Possui exatamente um kernel por núcleo
                        \item O estado é replicado entre os núcleos, não compartilhado
                        \item Comunicação explícita por mensagens entre monitores
                    \end{itemize}
                \end{column}
                \begin{column}{0.55\textwidth}
                    \centering
                    \begin{tikzpicture}[x=1mm,y=1mm]
                        % maquina
                        \draw[ufpr, thick, rounded corners=3pt] (-35,31) rectangle (35,-20);
                        % aplicacoes
                        \draw[fill=ufprlight, draw=kdraw, rounded corners=2pt] (-32,28) rectangle (32,22);
                        \node[font=\scriptsize\bfseries, text=ink] at (0,25) {Aplicações};
                        % nucleos
                        \foreach \x/\arch in {-25.5/x86,-8.5/x86,8.5/ARM,25.5/GPU} {
                            \draw[fill=ufpr!20, draw=kdraw, rounded corners=1pt] (\x-6,19) rectangle (\x+6,14);
                            \node[font=\tiny\bfseries] at (\x,16.5) {Réplica};
                            \draw[fill=white, draw=kdraw, rounded corners=1pt] (\x-6,12) rectangle (\x+6,6.5);
                            \node[font=\tiny\bfseries, text=ink] at (\x,9.25) {Monitor};
                            \draw[fill=kfill, draw=kdraw, rounded corners=1pt] (\x-6,1.5) rectangle (\x+6,-4);
                            \node[font=\tiny\bfseries] at (\x,-1.25) {CPU driver};
                            \draw[fill=hostc, draw=kdraw] (\x-6,-8) rectangle (\x+6,-13);
                            \node[font=\tiny\bfseries] at (\x,-10.5) {\arch};
                        }
                        % mensagens entre monitores
                        \foreach \a/\b in {-19.5/-14.5,-2.5/2.5,14.5/19.5} \draw[arr, <->] (\a,9.25) -- (\b,9.25);
                    \end{tikzpicture}
                \end{column}
            \end{columns}
        \end{frame}

% =====================================================================
        \section{Introdução}

        \begin{frame}{Arquitetura multikernel: multikernel no Linux}
            \begin{columns}[c]
                \begin{column}{0.50\textwidth}
                    \centering
                    \begin{tikzpicture}[x=6.4mm,y=6.4mm]
        \draw[ufpr, thick, rounded corners] (-0.3,-3.6) rectangle (10.3,3.0);
        \node[font=\scriptsize\bfseries, ufpr, anchor=west] at (-0.2,2.6) {Máquina Física};
        % nucleos
        \foreach \i in {0,...,9} {
            \ifnum\i<2 \def\c{hostc}\else\ifnum\i<5 \def\c{ufprlight}\else\def\c{ufpr!45}\fi\fi
            \fill[\c, draw=kdraw] (\i+0.08,0.7) rectangle (\i+0.92,1.5);
        }
        \draw[thick, gray, |-|] (0.08,1.75) -- (9.92,1.75);
        \node[font=\tiny, gray] at (5,2.1) {núcleos};
        % kernels
        \draw[kdraw, fill=hostc, rounded corners=2pt] (0.08,-0.6) rectangle (1.92,0.4);
        \node[font=\tiny\bfseries, align=center] at (1,-0.1) {Kernel\\hospedeiro};
        \draw[kdraw, fill=ufprlight, rounded corners=2pt] (2.08,-0.6) rectangle (4.92,0.4);
        \node[font=\tiny\bfseries, align=center] at (3.5,-0.1) {Kernel A\\(tempo real)};
        \draw[kdraw, fill=ufpr!45, rounded corners=2pt] (5.08,-0.6) rectangle (9.92,0.4);
        \node[font=\tiny\bfseries, align=center] at (7.5,-0.1) {Kernel B\\(propósito geral)};
        % memoria
        \fill[hostc, draw=kdraw] (0.08,-2.0) rectangle (1.92,-1.2);
        \fill[ufprlight, draw=kdraw] (2.08,-2.0) rectangle (4.92,-1.2);
        \fill[ufpr!45, draw=kdraw] (5.08,-2.0) rectangle (9.92,-1.2);
        \foreach \x in {1,3.5,7.5} \node[font=\tiny\bfseries] at (\x,-1.6) {memória};
        \draw[thick, gray, |-|] (2.08,-2.3) -- (9.92,-2.3);
        \node[font=\tiny, gray] at (6,-2.8) {instâncias: núcleos e memória dedicados};
                    \end{tikzpicture}
                \end{column}
                \begin{column}{0.50\textwidth}
                    \begin{itemize}
                        \item Quantidade configurável de núcleos dedicados com memória, imagem e kernel próprios
                        \item Kernels com configurações distintas e especializadas
                        \item Núcleos e memória realocáveis
                        \item Kernel de dispositivos centraliza E/S e dispositivos virtuais
                    \end{itemize}
                    \vspace{1mm}
                \end{column}
            \end{columns}
        \end{frame}

% =====================================================================
        \begin{frame}{Contêineres}
            \begin{columns}[c]
                \begin{column}{0.50\textwidth}
                    \small Empacotam a aplicação e suas dependências, executando como processos isolados sobre um mesmo kernel,
                                  utilizando \textit{namespaces} para isolamento e \textit{cgroups} para controlar o limite de recursos.

                    \vspace{2mm}
                    \setlength{\leftmargini}{1em}%
                    \textbf{Vantagens}
                    \begin{itemize}
                        \item Inicialização em milissegundos
                        \item Escalabilidade alta por nó
                    \end{itemize}
                    \vspace{1mm}
                    \textbf{Desvantagens}
                    \begin{itemize}
                        \item Interferência entre cargas
                        \item Impossibilidade de kernel especializado
                    \end{itemize}
                    \vspace{1mm}
                \end{column}
                \begin{column}{0.50\textwidth}
                    \centering
                    \begin{tikzpicture}[x=1.25mm,y=1.25mm]
                        \draw[ufpr, thick, rounded corners=3pt] (-25,3) rectangle (25,-40);
        \draw[fill=ufprlight, draw=kdraw, rounded corners=2pt] (-22,0) rectangle (-1,-9);
                        \node[chip] at (-11.5,-4.5) {App A + Dependências};
                        \draw[fill=ufprlight, draw=kdraw, rounded corners=2pt] (1,0) rectangle (22,-9);
                        \node[chip] at (11.5,-4.5) {App B + Dependências};
                        \draw[fill=kfill, draw=kdraw, rounded corners=2pt] (-22,-11) rectangle (22,-17);
                        \node[ltitle] at (0,-14) {Runtime (Docker)};
                        \draw[fill=kfill, draw=kdraw, rounded corners=2pt] (-22,-19) rectangle (22,-29);
                        \node[ltitle] at (0,-21.5) {Kernel do hospedeiro};
                        \node[tag] at (-9,-25.5) {namespaces}; \node[tag] at (9,-25.5) {cgroups};
                        \draw[fill=kfill, draw=kdraw, rounded corners=2pt] (-22,-31) rectangle (22,-37);
                        \node[ltitle] at (0,-34) {Hardware};
                    \end{tikzpicture}
                \end{column}
            \end{columns}
        \end{frame}

% =====================================================================
        \begin{frame}{Máquinas virtuais}
            \begin{columns}[c]
                \begin{column}{0.50\textwidth}
                    \centering
                    \begin{tikzpicture}[x=1.1mm,y=1.1mm]
                        \draw[ufpr, thick, rounded corners=3pt] (-25,3) rectangle (25,-37);
        \draw[fill=ufprlight, draw=kdraw, rounded corners=2pt] (-22,0) rectangle (-1,-7);
                        \node[chip] at (-11.5,-3.5) {App};
                        \draw[fill=ufprlight, draw=kdraw, rounded corners=2pt] (1,0) rectangle (22,-7);
                        \node[chip] at (11.5,-3.5) {App};
                        \draw[fill=kfill, draw=kdraw, rounded corners=2pt] (-22,-9) rectangle (-1,-16);
                        \node[chip] at (-11.5,-12.5) {Kernel convidado};
                        \draw[fill=kfill, draw=kdraw, rounded corners=2pt] (1,-9) rectangle (22,-16);
                        \node[chip] at (11.5,-12.5) {Kernel convidado};
                        \draw[fill=kfill, draw=kdraw, rounded corners=2pt] (-22,-18) rectangle (22,-25);
                        \node[ltitle] at (0,-21.5) {Hipervisor};
                        \draw[fill=kfill, draw=kdraw, rounded corners=2pt] (-22,-27) rectangle (22,-34);
                        \node[ltitle] at (0,-30.5) {Hardware};
                    \end{tikzpicture}
                \end{column}
                \begin{column}{0.50\textwidth}
                    \small Cada VM executa um sistema operacional completo, com seu próprio kernel, sobre um hipervisor
                    que virtualiza o hardware e intercepta instruções privilegiadas.

                    \vspace{2mm}
                    \setlength{\leftmargini}{1em}%
                    \textbf{Vantagens}
                    \begin{itemize}
                        \item Kernel próprio e configurável por VM
                        \item Isolamento forte entre VMs garantido pelo hipervisor
                    \end{itemize}
                    \vspace{1mm}
                    \textbf{Desvantagens}
                    \begin{itemize}
                        \item Sobrecarga de virtualização (instruções, E/S, rede, memória)
                        \item Boot lento, na ordem de segundos
                    \end{itemize}
                    \vspace{1mm}
                \end{column}
            \end{columns}
        \end{frame}

% =====================================================================
        \begin{frame}{Máquinas virtuais, contêineres e multikernel}
            \centering
            \begin{tikzpicture}[x=1mm,y=1mm, scale=0.86, every node/.style={transform shape}]
                \begin{scope}[shift={(-4,0)}]
                % ---------------- Maquinas virtuais
                \draw[card] (-67,0) rectangle (-23,-66);
                \node[head] at (-45,-6) {MÁQUINAS VIRTUAIS};
                \layer{-45}{-11}{10}{ufprlight}{kdraw}{Aplicações}
                \node[tag] at (-49,-17.5) {App}; \node[tag] at (-41,-17.5) {App};
                \layer{-45}{-23}{10}{kfill}{kdraw}{Kernels convidados}
                \node[tag] at (-50,-29.5) {Kernel}; \node[tag] at (-40,-29.5) {Kernel};
                \layer{-45}{-35}{10}{kfill}{kdraw}{Hipervisor + hospedeiro}
                \node[tag] at (-51,-41.5) {Traduções}; \node[tag] at (-38,-41.5) {E/S emulada};
                \layer{-45}{-47}{7}{kfill}{kdraw}{Hardware}
                \node[mkcap] at (-45,-60) {Dois kernels e um hipervisor\\entre a aplicação e o hardware};

                \end{scope}

                % ---------------- Conteineres
                \draw[card] (-22,0) rectangle (22,-66);
                \node[head] at (0,-6) {CONTÊINERES};
                \layer{0}{-11}{10}{ufprlight}{kdraw}{Aplicações}
                \node[tag] at (-4,-17.5) {App}; \node[tag] at (4,-17.5) {App};
                \layer{0}{-23}{10}{kfill}{kdraw}{Kernel compartilhado}
                \node[tag] at (0,-29.5) {Kernel único};
                \layer{0}{-35}{7}{kfill}{kdraw}{Hardware}
                \node[mkcap] at (0,-61) {Um kernel compartilhado\\por todas as cargas};

                \begin{scope}[shift={(4,0)}]
                % ---------------- Multikernel
                \draw[card] (23,0) rectangle (67,-66);
                \node[head] at (45,-6) {MULTIKERNEL};
                \layer{45}{-11}{10}{ufprlight}{kdraw}{Aplicações}
                \node[tag] at (41,-17.5) {App}; \node[tag] at (49,-17.5) {App};
                \layer{45}{-23}{16}{kfill}{kdraw}{Kernels independentes}
                \node[tag] at (39,-29.5) {Kernel A}; \node[tag] at (51,-29.5) {Kernel B};
                \node[tag] at (45,-35.5) {Kernel de dispositivos};
                \layer{45}{-41}{15}{kfill}{kdraw}{Hardware}
                \node[tag] at (45,-47.5) {Núcleos dedicados};
                \node[tag] at (45,-52.5) {Filas de hardware};
                \node[mkcap] at (45,-61) {Nada entre a aplicação\\e os núcleos};
                \end{scope}
            \end{tikzpicture}

            \vspace{-1mm}
            \refs{Multikernel Technologies, \url{https://multikernel.io/technology.html}}
        \end{frame}

% =====================================================================
        \begin{frame}{MicroVMs}
            \begin{columns}[c]
                \begin{column}{0.50\textwidth}
                    \centering
                    \begin{tikzpicture}[x=1mm,y=1mm]
                        \draw[ufpr, thick, rounded corners=3pt] (-25,3) rectangle (25,-39);
        \draw[fill=ufprlight, draw=kdraw, rounded corners=2pt] (-22,0) rectangle (-1,-7);
                        \node[chip] at (-11.5,-3.5) {App};
                        \draw[fill=ufprlight, draw=kdraw, rounded corners=2pt] (1,0) rectangle (22,-7);
                        \node[chip] at (11.5,-3.5) {App};
                        \draw[fill=kfill, draw=kdraw, rounded corners=2pt] (-22,-8.5) rectangle (-1,-15);
                        \node[chip] at (-11.5,-11.75) {Kernel mínimo};
                        \draw[fill=kfill, draw=kdraw, rounded corners=2pt] (1,-8.5) rectangle (22,-15);
                        \node[chip] at (11.5,-11.75) {Kernel mínimo};
                        \draw[fill=kfill, draw=kdraw, rounded corners=2pt] (-22,-16.5) rectangle (-1,-23);
                        \node[chip] at (-11.5,-19.75) {Hipervisor mínimo};
                        \draw[fill=kfill, draw=kdraw, rounded corners=2pt] (1,-16.5) rectangle (22,-23);
                        \node[chip] at (11.5,-19.75) {Hipervisor mínimo};
                        \draw[fill=kfill, draw=kdraw, rounded corners=2pt] (-22,-24.5) rectangle (22,-30);
                        \node[ltitle] at (0,-27.25) {Kernel hospedeiro};
                        \draw[fill=kfill, draw=kdraw, rounded corners=2pt] (-22,-31.5) rectangle (22,-36);
                        \node[ltitle] at (0,-33.75) {Hardware};
                    \end{tikzpicture}
                \end{column}
                \begin{column}{0.50\textwidth}
                    \begin{itemize}
                        \item MicroVMs são instâncias mínimas, possuindo apenas dispositivos essenciais para a aplicação
                        \item Referência para avaliação experimental: kernel próprio por unidade e orquestração madura
                        \item Comparadas em condições equivalentes, mesmos núcleos, memória, nós e cargas
                    \end{itemize}
                \end{column}
            \end{columns}
        \end{frame}

% =====================================================================
        \begin{frame}{Orquestração e o problema}
            \begin{columns}[T]
                \begin{column}{0.50\textwidth}
                    \centering
                    \begin{itemize}
                        \item Contêineres e VMs ganharam muita popularidade por conta dos orquestradores
                        \item Orquestradores focam em manter um estado desejado e balancear suas unidades, além de monitorá-las e
                        prover meios de escalar suas instâncias
                        \item As pesquisas em multikernel sempre tiveram foco em uma única máquina e na arquitetura em si
                    \end{itemize}
                \end{column}
                \begin{column}{0.50\textwidth}
                    \centering
                    \begin{tikzpicture}[x=1mm,y=1mm]
                        \tikzset{tool/.style={dbox, minimum width=24mm, minimum height=6mm, font=\scriptsize}}
                        \draw[ufpr, thick, rounded corners=3pt] (-28,6) rectangle (2,-30);
                        \node[tool] at (-13,0) {Kubernetes};
                        \node[tool] at (-13,-8) {Docker Swarm};
                        \node[tool] at (-13,-16) {Apache Mesos};
                        \node[tool] at (-13,-24) {...};
                    \end{tikzpicture}
                    \vspace{2mm}
                \end{column}
            \end{columns}
        \end{frame}

% =====================================================================
        \begin{frame}{Pergunta de pesquisa e hipótese}
            \small
            \textbf{Pergunta} -- Quais mecanismos são necessários para provisionar, posicionar e monitorar instâncias
            multikernel em vários nós,
            e qual o custo desse gerenciamento em relação à orquestração de microVMs?\\
            \vspace{10mm}
            \textbf{Hipótese} -- É possível orquestrar instâncias multikernel em vários nós de modo que:
            \begin{itemize}
                \item O custo introduzido pelo orquestrador seja uma pequena fração do tempo total de provisionamento (H1)
                \item As instâncias apresentem menos interferência de desempenho entre si do que microVMs, sob as mesmas
                cargas e orquestração equivalente (H2)
                \item O particionamento dinâmico aumente a utilização dos nós em relação ao estático (H3)
            \end{itemize}
        \end{frame}

        \section{Proposta}

        \begin{frame}{Proposta: K-ORCH}
            \begin{columns}[T]
                \begin{column}{0.50\textwidth}
                    \textbf{Funcionamento:}
                    \setlength{\leftmargini}{1em}
                    \begin{itemize}
                        \item Pedido descrevendo o estado desejado
                        \item Escolha do nó por núcleos e memória livres, e por critérios de balanceamento
                        \item Agente particiona o hardware local e inicializa a instância
                        \item Detecção de falhas e recuperação
                        \item Escolha de particionamento estático ou dinâmico
                    \end{itemize}
                \end{column}
                \begin{column}{0.50\textwidth}
                    \textbf{Escopo:}
                    \setlength{\leftmargini}{1em}
                    \begin{itemize}
                        \item Orquestração em uma rede local, sob domínio administrativo único
                        \item Fora do escopo: \textit{datacenters}, múltiplos domínios e comunicação pela Internet
                        \item Isolamento e segurança entre instâncias pertencem ao modelo da arquitetura e não ao escopo da pesquisa
                    \end{itemize}
                \end{column}
            \end{columns}
        \end{frame}

% =====================================================================
        \begin{frame}{Objetivos}
            \begin{itemize}
                \item Levantar o estado da arte
                \item Definir formalmente a unidade orquestrável
                \item Prototipar a arquitetura
                \item Implementar o ciclo de vida e o posicionamento
                \item Implementar o particionamento
                \item Implementar o monitoramento
                \item Avaliar experimentalmente
            \end{itemize}
        \end{frame}

% =====================================================================
        \section{Metodologia}

        \begin{frame}{Metodologia}
            \framesubtitle{Pesquisa teórico-prática, quantitativa e exploratória}
            \begin{columns}[T,onlytextwidth]
                \begin{column}{0.57\textwidth}
                    \setlength{\leftmargini}{1em}%
                    \textbf{1. Caracterizar}
                    \begin{itemize}
                        \item Estado da arte: orquestradores e gerência de multikernels
                        \item Multikernel com versão e hardware fixados
                    \end{itemize}
                    \vspace{1mm}
                    \textbf{2. Prototipar}
                    \begin{itemize}
                        \item Unidade orquestrável: núcleos, memória, imagem, kernel e estado desejado
                        \item Arquitetura: plano de controle, agentes e comunicação
                        \item Ciclo de vida (criar, iniciar, parar e remover) e posicionamento \textit{first-fit} ou duas escolhas
                    \end{itemize}
                \end{column}
                \begin{column}{0.39\textwidth}
                    \setlength{\leftmargini}{1em}%
                    \textbf{3. Particionar e monitorar}
                    \begin{itemize}
                        \item Particionamento estático ou dinâmico
                        \item Detecção de falhas: reiniciar ou reposicionar a instância
                    \end{itemize}
                    \vspace{1mm}
                    \textbf{4. Avaliar}
                    \begin{itemize}
                        \item Experimentos em nós físicos
                        \item Comparar a microVMs e entre modos de particionamento
                    \end{itemize}
                \end{column}
            \end{columns}
        \end{frame}

% =====================================================================
        \begin{frame}{Avaliação experimental}
            \centering
            \vspace{2mm}
            \renewcommand{\arraystretch}{1.9}
            \small
            \begin{tabular}{>{\raggedright\arraybackslash}p{2.5cm}>{\raggedright\arraybackslash}p{5.0cm}>{\raggedright\arraybackslash}p{4.0cm}}
                \textbf{Experimento} & \textbf{Métrica} & \textbf{Avaliação} \\
                \hline
                Provisionamento           & tempo do pedido até a instância estar saudável  & custo do orquestrador (H1) \\
                Densidade e interferência & instâncias por nó e latência sob grandes cargas & isolamento de desempenho (H2) \\
                Particionamento           & utilização dos nós e pedidos rejeitados         & estático $\times$ dinâmico (H3) \\
                Falhas                    & detecção, recuperação e taxa de sucesso         & recuperação \\
                \hline
            \end{tabular}
        \end{frame}

% =====================================================================
        \section{Resultados esperados e cronograma}

        \begin{frame}{Resultados esperados}
            \begin{itemize}
                \item Protótipo do K-ORCH, com código e scripts públicos
                \item Modelo de instância e de estado desejado
                \item Custos e limites da orquestração em comparação com microVMs
                \item Divulgação em artigo para evento ou periódico
                \item Ganhos esperados: menor interferência entre instâncias e kernels especializados por carga, com maior desempenho
                \item Custos esperados: menor densidade por nó e tempo de boot maior
            \end{itemize}
        \end{frame}

% =====================================================================
        \begin{frame}{Cronograma}
            \centering
            \scriptsize
            \renewcommand{\arraystretch}{1.12}
            \setlength{\tabcolsep}{3.5pt}
            \begin{tabular}{p{8.4cm}cccccccc}
                \hline
                \textbf{Atividade} & T1 & T2 & T3 & T4 & T5 & T6 & T7 & T8 \\
                \hline
                Levantamento do estado da arte & $\bullet$ & $\bullet$ &  &  &  &  &  &  \\
                Caracterização experimental da plataforma multikernel & $\bullet$ & $\bullet$ &  &  &  &  &  &  \\
                Definição da unidade orquestrável e do modelo de estado &  & $\bullet$ &  &  &  &  &  &  \\
                Projeto da arquitetura &  &  & $\bullet$ &  &  &  &  &  \\
                Implementação do ciclo de vida e do posicionamento &  &  & $\bullet$ & $\bullet$ &  &  &  &  \\
                Redação, entrega e defesa do TCC &  &  & $\bullet$ & $\bullet$ &  &  &  &  \\
                Implementação do particionamento de recursos &  &  &  &  & $\bullet$ & $\bullet$ &  &  \\
                Implementação do monitoramento e da recuperação de instâncias &  &  &  &  & $\bullet$ & $\bullet$ &  &  \\
                Preparação do ambiente de avaliação e das microVMs de referência &  &  &  &  & $\bullet$ & $\bullet$ &  &  \\
                Avaliação experimental e comparativa &  &  &  &  &  &  & $\bullet$ & $\bullet$ \\
                Refinamento do orquestrador &  &  &  &  &  &  & $\bullet$ & $\bullet$ \\
                Redação e submissão de artigos &  &  &  &  &  & $\bullet$ & $\bullet$ & $\bullet$ \\
                Defesa do mestrado &  &  &  &  &  &  &  & $\bullet$ \\
                \hline
            \end{tabular}
        \end{frame}



    }

% =====================================================================
    {%
        \usebackgroundtemplate{\centering\includegraphics[width=\paperwidth]{CapaUFPR.png}}
        \setbeamertemplate{footline}{}
        \setbeamertemplate{headline}{\vskip2mm\hfill{\footnotesize\color{gray}\insertframenumber\,/\,\inserttotalframenumber}\hspace*{4mm}}
        \begin{frame}
            \centering
            \vspace{-14mm}
            {\Huge\bfseries\color{ufpr} Obrigado!}\\[4mm]
            {\small\textbf{Eduardo Kaluf}}\\[1mm]
            {\small eduardokaluf@ufpr.br}\\[1mm]
            {\footnotesize Orientador: Prof. Carlos A. Maziero}
        \end{frame}
    }

\end{document}